% 本文件由 LaTeX 排版 Agent 根据有效 translation.json 自动生成。
% author=Apocrypha; work_id=0848; title=伪玛窦福音

\chapter{伪玛窦福音}

% structure: p0001 source=h1 target=omitted reason=page-title-already-rendered-by-parent

《真福玛利亚诞生及救主童年史》开始。由真福福音作者玛窦用希伯来文写成，并由真福司铎热罗尼莫翻译为拉丁文。\par

致他们亲爱的弟兄热罗尼莫司铎，克罗玛提乌斯主教与赫略多鲁斯主教在主内问安。\par

我们在伪经中找到了童贞玛利亚的诞生及我们主耶稣基督的诞生与童年。但考虑到其中写有许多违反我们信仰的事，我们相信应当全部拒绝，免得我们将基督的喜乐转给敌基督。因此，当我们考虑这些事时，有圣者帕尔梅尼乌斯和瓦里努斯前来，说您的圣德找到了一卷希伯来文手稿，由最真福的宗徒玛窦亲手所写，其中也记载了童贞母亲本人的诞生和我们救主的童年。因此，我们藉我们主耶稣基督亲自恳求您的爱德，将其从希伯来文译成拉丁文，与其说是为了获取那些属于基督标记的事，不如说是为了排除异端的诡计，他们为教坏道理，将自己编造的谎言与基督卓越的诞生混杂，企图以生命的甘甜隐藏死亡的苦涩。因此，您最纯洁的虔诚理当倾听我们这些弟兄的恳求，或让我们这些主教要求您所认为应尽的爱德之债。\par

\Emph{热罗尼莫给他们的回信。}\par

致我圣洁至福的主教克罗玛提乌斯与赫略多鲁斯，热罗尼莫，基督的卑微仆人，在主内问安。\par

一个知道地下有金矿的人，不会立刻攫取翻起的沟壑所涌出的一切；但在颤抖的锄头举起闪光的金块之前，他暂且停留，翻动并举起土块，尤其是在他并未增益收益之时。一项艰巨的任务加在我身上，因为您的真福所命令我的，圣宗徒与福音作者玛窦本人并未为出版而写。因为如果他没有多少秘密地这样做，他也会将其加在他所出版的福音中。但他用希伯来文写了这卷书，而出版得如此之少，以至于今日这卷由他亲手用希伯来文所写的书，仍保存在非常虔诚的人手中，代代相传。他们从未将此书交给任何人翻译。结果，当它由一位名叫路济乌斯的摩尼门徒出版时——此人亦写了那部伪作的《宗徒大事录》——这卷书提供的不是建树，而是毁灭；而主教会议对它的意见，按其应得，是教会不应为其敞开耳朵。现在，那些对我们狂吠之人的咆哮当止息；因为我们并非将这卷小书加于正典著作，而是翻译一位宗徒与福音作者所写，为揭露异端的谎言。在此工作中，我们既服从虔诚主教的命令，又对抗不虔敬的异端。因此，我们所实践的是基督的爱德，相信那些因我们的顺从而得以认识我们救主神圣童年的人，必以祈祷帮助我们。\par

另有一封致同一主教的信，相传出自热罗尼莫：——\par

您问我让您知道我对一本被一些人认为是关于圣玛利亚诞生的书的看法。因此，我希望您知道其中有许多虚假。因为一位名叫塞琉古的人——他写了《宗徒的苦难》——写了这本。但是，正如他关于他们的能力和所行的奇迹写了真实之事，但关于他们的教义说了许多虚假；同样，在这里他也凭空捏造了许多不实之言。我将小心地逐字翻译，正如它在希伯来文中一样，因为据称它是由圣福音作者玛窦所写，以希伯来文写成，并置于他的福音之前。这是真是假，我留给序言作者和作者的可靠性；至于我自己，我宣布它们可疑；我不断言它们明显虚假。但我坦然说——我认为没有信友会否认——无论这些故事是真是假，圣玛利亚的神圣诞生之前有伟大奇迹，之后有更大奇迹；因此，那些相信天主能做这些事的人，可以相信并阅读它们，而不损害他们的信德或危及他们的灵魂。简言之，就我所能，宁追作者之意而非其辞，有时行走同路而无同一足迹，有时稍离正道，但仍保持同一途径，我将以此行文叙事，不会说任何那里未写、或按同一思路可能写下的事。\par

% structure: p0010 source=h2 target=section reason=relative-heading-rank; base=h2

\section{第一章}

在那些日子，耶路撒冷有一个人，名叫约阿基姆，属犹大支派。他牧养自己的羊群，以纯全而专一的心敬畏主。他除了自己的牲畜之外别无所虑，用它们的出产供给所有敬畏天主的人吃喝，在敬畏天主中将双倍礼物献给所有在教义上劳苦并事奉祂的人。因此，他的羔羊、他的羊、他的羊毛，以及他所有的一切，他都分作三份：一份给孤儿、寡妇、旅客和穷人；第二份给敬拜天主的人；第三份留给自己和全家。\BibleRef{多 1:7} 他这样行，主就使他的牲畜增多，以致在以色列民中无人像他。他十五岁时开始这样做。二十岁时娶了同支派亚甲的女儿亚纳，即犹大支派、达味家族。他们同居二十年，他却无儿无女。\par

% structure: p0012 source=h2 target=section reason=relative-heading-rank; base=h2

\section{第二章}

正值庆节期间，在那些向主献香的人中，\Person{约雅敬}站在主前预备自己的礼品。有位名叫\Person{如本}的司铎来到他跟前说：\Quote{你不可以站在那些向天主献祭的人中，因为天主没有降福你使你在以色列中有子嗣。}于是他在众人面前受辱，哭着离开主的圣殿，没有回家，而是带着牧人往远方的山区自己的羊群那里去，以致他的妻子\Person{亚纳}五个月没有他的音讯。她流着泪祈祷说：\Quote{主，以色列的全能天主，既然你还没有赐给我儿女，为什么又把我的丈夫夺去呢？看，我已经五个月没有见到我的丈夫了；我不知道他在哪里停留；即使我知道他死了，也不能埋葬他。}她痛哭不止，进入圣殿的庭院，俯伏在地祈祷，在主前倾心吐意。之后，她从祈祷中起来，举目向天主，看见月桂树上有一个麻雀窝\BibleRef{多 2:10}，就叹息着向主扬声说：\Quote{上主全能的天主，你赐给了所有受造物子嗣，无论野兽家畜、蛇、鸟和鱼，它们都为幼崽欢欣，唯独你将我摒绝于你仁慈的恩赐之外。因为天主，你知道我的心，从我婚姻之初，我就许愿说，如果天主赐给我儿子或女儿，我必将他们奉献于你的圣殿。}她正这样说时，忽然主的使者出现在她面前说：\Quote{\Person{亚纳}，不要害怕，因为在天主的旨意中你有子嗣；所有世代直到终结，都要因你将要生的子而感到惊异。}他说完这话，就从她眼前消失了。但她因见了这样的景象、听了这样的话而恐惧战栗，最后进入卧室，倒在床上如同死人。她一整天一整夜都在极大的战栗和祈祷中。之后，她叫来自己的婢女，对她说：\Quote{你看见我寡居受骗、极度困惑，却不愿进来见我？}那婢女低声喃喃回答说：\Quote{如果天主封闭了你的子宫，又夺去了你的丈夫，我能为你做什么呢？}\Person{亚纳}听了这话，就放声大哭。\par

% structure: p0014 source=h2 target=section reason=relative-heading-rank; base=h2

\section{第三章}

与此同时，在山上有一位青年出现在正在放羊的\Person{约雅敬}面前，对他说：\Quote{你为什么不回到你的妻子那里去？}\Person{约雅敬}说：\Quote{我娶她已经二十年，天主没有愿意我藉她得子嗣。我带着羞辱和责骂被赶出主的圣殿：既然我已一度被抛弃、被完全藐视，为什么还要回到她那里去？我就在这里与我的羊群同住；只要天主愿意赐我在世的光明，我愿意藉我仆人的手，将我的份给予穷人、孤儿和敬畏天主的人。}他这样说时，那青年对他说：\Quote{我是主的使者，我今日已在你妻子哭泣祈祷时显现给她，安慰了她；你要知道，她已从你的后裔怀了一个女儿，你却因不知情而离开了她。她将住在天主的圣殿，圣神将住在她的内；她的福气将大于所有圣善的妇女，以致没有人能说在她之前有像她的，或在她之后在这世上会有如此者。因此，你下山回到你的妻子那里去，你将发现她已怀孕。因为天主已在她的身上兴起了后裔，为此你要感谢天主；她的后裔必受祝福，她自己也必受祝福，并要成为永恒祝福之母。}于是\Person{约雅敬}朝拜了天使，对他说：\Quote{如果我在你眼中蒙恩，请在我的帐幕里坐一会儿，祝福你的仆人。}\BibleRef{创 18:3}天使对他说：\Quote{不要说仆人，而是同仆；因为我们都是同一主人的仆人。\BibleRef{默 19:10}但我的食物是看不见的，我的饮品是凡人所不能见的。因此你不该请我进入你的帐幕；但如果你要给我什么东西，\BibleRef{民 13:16}就当作全燔祭献给主。}于是\Person{约雅敬}取了一只无瑕疵的羔羊，对天使说：\Quote{我不敢向主献全燔祭，除非你的命令赐予我献祭的司祭权。}天使对他说：\Quote{我不会邀请你献祭，除非我知道主的旨意。}当\Person{约雅敬}向天主献祭时，天使和祭品的馨香一同随着烟气直上天堂。\BibleRef{民 13:20}\par

于是\Person{约雅敬}俯伏在地，从白天的第六时辰直到傍晚躺卧祈祷。与他同来的少年和雇工看见他，不知他为何躺卧，以为他死了；他们来到他跟前，好不容易才把他从地上扶起来。他向他们讲述了天使的神视，他们大为恐惧和惊奇，劝他立即完成天使的神视，急速回到他的妻子那里去。当\Person{约雅敬}心中犹豫是否回去时，他沉沉入睡；看，那在他醒时曾向他显现的天使，在梦中向他显现说：\Quote{我是天主任命守护你的天使：放心下山，回到\Person{亚纳}那里去，因为你和你妻子\Person{亚纳}所行的仁爱善工已在至高者面前被述说；天主将赐给你这样的果实，是从起初没有先知或圣人曾拥有过的，将来也不会有。}\Person{约雅敬}从梦中醒来，把所有的牧人都叫来，告诉他们他的梦。他们便朝拜主，对他说：\Quote{你要注意，不再轻看天使的话。但起来，让我们离开这里，安静地回去，放牧我们的羊群。}\par

当他们用三十天的时间返回，如今临近家门时，看，上主的天使出现在正站着祈祷的亚纳面前，说：\BibleRef{宗 9:11}\Quote{你到那称为金门的门口去，在路上迎接你的丈夫，因为他今日要你这里来。}于是她急忙带着侍女向他走去，并向上主祈祷，在门口站了很久等候他。当她因久候而疲惫时，举目看见约阿基姆远远地带着羊群走来；她便跑向他，搂住他的脖子，感谢天主说：\Quote{我曾是寡妇，看，如今不再是了；我曾是不孕的，看，如今我已怀了孕。}于是他们朝拜了主，回到自己家中。这事传开后，所有的邻人和相识都大为喜乐，以致全以色列地都为他们庆贺。\par

% structure: p0018 source=h2 target=section reason=relative-heading-rank; base=h2

\section{第四章}

此后，她的孕期满了九个月，亚纳生了一个女儿，给她取名叫玛利亚。在她三岁断奶后，约阿基姆和他的妻子亚纳一同上主的圣殿去，向天主献上祭献，并将这婴孩——名叫玛利亚——安置在童贞女团体内，这些童贞女日日夜夜留在那里赞美天主。当她被放在圣殿门前时，她如此迅速地登上了十五级台阶，丝毫没有回头，也没有像孩童常做的那样寻找父母。于是她的父母都焦急地寻找孩子，两人同样惊讶，直到在圣殿里找到了她，连圣殿的司铎们也感到惊奇。\par

% structure: p0020 source=h2 target=section reason=relative-heading-rank; base=h2

\section{第五章}

那时，亚纳充满圣神，在众人面前说：\Quote{全能的上主，万军的天主，记念祂的圣言，以良善而神圣的眷顾眷顾了祂的子民，为要挫败那起来攻击我们的外邦人的心，使他们归向祂。祂垂听了我们的祈祷；祂使我们众仇敌的夸耀远离了我们。不孕的成了母亲，为以色列带来了欢跃和喜乐。看，这些就是我带来献给我的主的礼物，我的仇敌未能阻止我。因为天主使他们的心转向了我，祂亲自赐给了我永久的喜乐。}\par

% structure: p0022 source=h2 target=section reason=relative-heading-rank; base=h2

\section{第六章}

玛利亚受到所有以色列人的赞叹；当她三岁时，她的步伐如此成熟，言语如此完美，且如此勤奋地致力于赞美天主，以致众人都对她感到惊奇；她不被人看作幼童，而如同一个三十岁的成年人。她在祈祷中如此恒心，她的容貌如此美丽而荣光，以至于几乎无人能直视她的面容。她经常从事羊毛纺织，以致在她幼年时就能做出老年妇女所不能做的事。她为自己立下了这样的规矩：从早晨到第三时辰，她留在祈祷中；从第三时辰到第九时辰，她从事纺织；从第九时辰起，她又致力于祈祷。她从不停止祈祷，直到上主的天使向她显现，她常从天使手中领受食物；因此她在天主的工程上日益成全。当较年长的童贞女停止赞美天主时，她绝不停止；因此在赞美天主和守夜里，无人比她更在前，无人比她更精通天主法律的智慧，更谦卑，更善于歌唱，更成全于一切美德。她确实坚定、不动摇、不变，日日进步于成全。无人见她发怒，也无人听她说恶言。她的一切言语如此充满恩宠，以致她的天主被人认出在其舌上。她常常祈祷并研读法律，且担心因自己的任何言语而在同伴身上犯罪。她又担心自己在笑声或美丽的声音中犯任何过失，或因得意而向同辈中任何人显露恶行或傲慢。她无间断地赞美天主；恐怕即使在问候时也中止赞美天主；若有人问候她，她常以问候的方式回答：\Quote{感谢天主。}从此，人们在互相问候时首先开始说\Quote{感谢天主}。她只以每日从天使手中领受的食物维持生活；但她从司祭那里得到的食物，则分给穷人。天主的使者常被看见与她交谈，并极其勤勉地服从她。若任何生病的人触摸了她，那人当时就痊愈回家了。\par

% structure: p0024 source=h2 target=section reason=relative-heading-rank; base=h2

\section{第七章}

那时，司铎阿彼雅塔尔向大司祭们不断献礼，为将玛利亚娶作他儿子的妻子。但玛利亚禁止他们，说：\Quote{我认识一个男人，或一个男人认识我，是不可能的。}因为所有的司铎和所有她的亲属不断对她说：\Quote{天主在子女中受敬拜，在后代中受钦崇，正如在以色列子民中一向如此。}但玛利亚回答他们说：\Quote{天主在贞洁中受敬拜，这是首先得到证明的。因为在亚伯尔之前，人中没有义人，他凭自己的祭品悦乐了天主，却被那使天主不悦的人无情地杀害了。因此他得了两顶冠冕——祭献的冠冕和童贞的冠冕，因为他的肉身没有玷污。厄里亚尚在肉身时，也被提升到肉身上去，因为他保存了自己的肉身无瑕。我自幼在上主的圣殿中便已学到，童贞足以令天主喜悦。因此，既然我能够献上令天主喜悦的，我心中决定——我绝不认识任何男人。}\par

% structure: p0026 source=h2 target=section reason=relative-heading-rank; base=h2

\section{第八章}

那时，玛利亚已十四岁，为此法利塞人有了藉口，说按惯例，那个年纪的女子不应再住在天主的殿里。于是他们决定派一个使者走遍以色列各支派，叫众人在第三天都聚集到上主的殿里。众人都聚集后，大司铎阿彼雅塔尔站起来，登上更高的一级台阶，好让众人看见并听见他。待全场肃静后，他说：\Quote{以色列子民，请听我说，把我的话收入你们耳中。自从撒罗满建造这殿以来，殿中一直有贞女，有君王的女儿，有先知的女儿，也有大司铎和司铎的女儿；她们都是伟大的，值得钦佩的。但当她们到了适当的年龄，就被嫁出去，仿效她们先前的母亲的行径，中悦天主。但唯有玛利亚找到了一种新的生活方式，她许愿要终身守贞事奉天主。因此在我看来，我们应当借着询问和天主的答复，设法弄清她应托付给谁看管。} 这话得到了全会众的赞同。司铎们便为十二支派抽签，签落在犹大支派。司铎说：\Quote{明天凡没有妻子的都来，手里拿着自己的棍杖。} 于是若瑟也同青年们一起带来了他的棍杖。棍杖都交给了大司铎，他向 上主天主 献了祭，并求问了上主。上主对他说：\Quote{把他们的棍杖都放在天主的至圣所里，让它们留在那里，吩咐他们明天来取回自己的棍杖；那人的棍杖顶端会飞出一只鸽子，飞向天空，并且当棍杖交还时，在他手里显出这个记号，就把玛利亚交给他看管。}\par

第二天，众人一早聚集，献了香祭后，大司铎进入至圣所，取出了棍杖。他把棍杖分发完毕，没有一根有鸽子飞出，大司铎就穿上十二个铃铛的祭服和司铎长袍，进入至圣所，在那里献了全燔祭，倾心祈祷。上主的天使显现给他，说：\Quote{这里有一根最短的棍杖，你没有理会；你同其他棍杖一起带了进来，却没有同它们一起取出去。你若把它取出来，交给它的主人，它里面就会显出我对你说过的记号。} 那根棍杖正是若瑟的；因为他年纪老了，似乎已经被弃绝，免得他接收玛利亚，而且他自己也不愿索回棍杖。他谦卑地站在最后，大司铎大声向他喊道：\Quote{来吧，若瑟，收回你的棍杖，因为我们正等你。} 若瑟战战兢兢地走上前来，因为大司铎用极大的声音叫他。他刚一伸手抓住自己的棍杖，立刻从棍杖顶端飞出一只鸽子，洁白如雪，极其美丽，在殿顶上下飞翔良久，终于飞向天空。于是众人都向老人道贺，说：\Quote{你年老有福了，若瑟父亲，因为天主显示你堪当接收玛利亚。} 司铎们对他说：\Quote{你接收她吧，因为在犹大支派中，唯独你被天主拣选了。} 若瑟开始羞怯地对他们说：\Quote{我是个老人，又有儿女；你们为什么把这个比我孙子还小的幼女交给我呢？} 大司铎阿彼雅塔尔对他说：\Quote{若瑟，你要记住，达堂、阿彼兰和科辣黑因轻视天主的旨意而灭亡。你若轻视天主的命令，也会如此。} 若瑟回答说：\Quote{我并没有轻视天主的旨意；但我要作她的监护人，直到我明白天主的旨意，看我的哪个儿子可以娶她为妻。请给她一些同伴贞女，让她暂且与她们同住，作为安慰。} 大司铎阿彼雅塔尔回答说：\Quote{的确，将给她五个贞女作为安慰，直到指定的日子你可以接收她；因为不能许配给其他人。}\par

于是若瑟接收了玛利亚，连同其他五个要与她同住在若瑟家的贞女。这些贞女是：黎贝加、塞福辣、苏撒纳、阿彼盖耳和加厄耳。大司铎给了她们丝绸、蓝色、细麻、朱红、紫色和细亚麻。因为她们抽签决定每个贞女做什么，为上主圣殿的帐幔织紫色布匹的签落在玛利亚身上。她拿到后，那些贞女对她说：\Quote{既然你最小、最卑微、比我们都年轻，你就配得获取和拥有紫色。} 她们这样说，仿佛带着恼怒的话语，开始称她为贞女之后。然而，正当她们这样做时，上主的天使出现在她们中间，说：\Quote{这些话不是出于恼怒说出的，而是作为最真实的预言预言的。} 因此，她们见了天使和他的话就战栗，请求玛利亚宽恕她们，并为她们祈祷。\par

% structure: p0030 source=h2 target=section reason=relative-heading-rank; base=h2

\section{第九章。}

第二天，玛利亚在水泉边打水时，上主的天使显现给她，说：\Quote{玛利亚，你是有福的，因为你在你的胎中为上主预备了居所。因为，看，天上的光将要来住在你内，并藉着你在全世界照耀。}\par

又有一天，当她用手指在紫色线上工作时，有一位难以言喻的美少年进来。\Person{玛利亚}看见他，就极其害怕并颤抖。他对她说：\Quote{万福，\Person{玛利亚}，充满恩宠，主与你同在；你在女人中应受赞颂，你胎中的果实也应受赞颂。}\BibleRef{路 1:28}她听了这些话，就颤抖，非常害怕。然后主的天使又说：\Quote{玛利亚，不要害怕；因为你在天主面前获得了宠幸。看，你将在胎中怀孕，并生一位君王，祂不仅充满大地，也充满天，且万代为王。}\par

% structure: p0033 source=h2 target=section reason=relative-heading-rank; base=h2

\section{第十章。}

当这些事发生时，\Person{若瑟}正在沿海地区的建筑工地上忙于他的房子建造工作；因为他是一个木匠。九个月后他回到家，发现\Person{玛利亚}怀孕了。因此，他极其痛苦，颤抖着喊道：\Quote{主天主，请收去我的灵魂；因为我与其苟活，不如死去。}与他在一起的童贞女们对他说：\Quote{若瑟，你在说什么？我们知道没有人碰过她；我们可以作证她仍是童贞女，未被玷污。我们一直守护她；她常同我们祈祷；每天天主的使者对她说话；每天她从主手中接受食物。我们不知道她怎么可能有任何罪。但如果你想要我们告诉你我们所怀疑的，那就是除主的天使外，无人使她怀孕。}于是\Person{若瑟}说：\Quote{你们为何误导我，让我相信主的天使使她怀孕？但可能有人假装成主的天使，欺骗了她。}说着，他哭了，说：\Quote{我有什么脸面仰望主的圣殿？有什么脸面见天主的司祭？我该怎么办？}说着，他想逃走，休了她。\par

% structure: p0035 source=h2 target=section reason=relative-heading-rank; base=h2

\section{第十一章。}

当他正在想站起来躲藏，并秘密居住时，看，当夜，主的天使在梦中向他显现，说：\Quote{若瑟，达味之子，不要害怕；你要娶\Person{玛利亚}为你的妻子；因为她腹中的是出于\OriginalTerm{圣神}{Holy Spirit}。她将生一个儿子，祂的名字要叫\Person{耶稣}，因为祂要把祂的人民从他们的罪恶中拯救出来。}\Person{若瑟}从梦中醒来，感谢天主，并与\Person{玛利亚}和与她在一起的童贞女们说话，告诉了他们他的异象。于是他因\Person{玛利亚}而得了安慰，说：\Quote{我犯了罪，因为我曾怀疑你。}\par

% structure: p0037 source=h2 target=section reason=relative-heading-rank; base=h2

\section{第十二章。}

此后，有传言说\Person{玛利亚}怀了孕。\Person{若瑟}被圣殿的官员抓住，同\Person{玛利亚}一起被带到\OriginalTerm{大司祭}{high priest}那里。他与司祭们开始责备他，说：\Quote{你为何欺骗如此伟大荣耀的童贞女？她像鸽子一样在圣殿被天主的使者喂养，从未想要看见或拥有男人，她对天主的律法有最卓越的知识。若不是你向她施行强暴，她本可保持童贞。}\Person{若瑟}发誓，他从未碰过她。\Person{阿比亚塔尔}大司祭回答他：\Quote{主永在，我要给你喝\OriginalTerm{上主饮用之水}{water of drinking of the Lord}，你的罪就立刻显现。}\par

于是聚集了数不清的群众，\Person{玛利亚}被带到圣殿。司祭、她的亲属和她的父母都哭了，对\Person{玛利亚}说：\Quote{向司祭告明你的罪吧，你曾像鸽子一样在天主的圣殿里，从天使手中接受食物。}\Person{若瑟}又被召到祭台前，给他喝\OriginalTerm{上主饮用之水}{water of drinking of the Lord}。凡撒谎者喝了这水，并绕祭台行走七圈，天主就会在他脸上显出某种记号。因此，当\Person{若瑟}安全地喝了水，并绕祭台七圈后，他身上没有显出任何罪的记号。于是所有司祭、官员和民众都称赞他无罪，说：\Quote{你有福了，因为你没有查出任何罪证。}他们又叫来\Person{玛利亚}，说：\Quote{你还有什么借口？或者还有什么比你的怀孕更能暴露你的记号呢？我们只要求你这一点，既然\Person{若瑟}在你这件事上是清白的，你就承认是谁欺骗了你。因为你的告白暴露你自己，总比天主的义怒在你脸上留下印记，并在众人面前揭露你更好。}然后\Person{玛利亚}坚定而无畏地说：\Quote{主天主，万有的君王，知晓一切秘密，如果我身上有任何玷污、任何罪、任何邪欲或不贞，请在众人面前揭露我，使我成为惩罚所有人的榜样。}说完，她勇敢地走上主的祭台，喝下那水，并绕祭台七圈，她身上没有找到任何瑕疵。\par

当所有的人都极其惊讶，看见她怀了孕，脸上却没有显现任何记号，他们开始因相互矛盾的看法而困惑：有的说她是圣洁无玷的，有的说她是邪恶污秽的。\Person{玛利亚}见自己仍被民众怀疑，而且因此似乎未能完全清白，就在众人面前大声说：\Quote{\OriginalTerm{阿多奈}{Adonai}上主永在，我站在祂面前的\OriginalTerm{万军的上主}{Lord of Hosts}，我不认识男人；但我自幼献身于祂，祂认识我。我从小就对天主立下这誓愿：在我创造者内保持无玷，我信靠我只为祂而活，只事奉祂；在我有生之年，我愿在祂内保持纯净。}于是众人开始吻她的脚，抱她的膝，求她宽恕他们邪恶的猜疑。她被民众、司祭和所有童贞女欢欣鼓舞地领回她的家。他们喊道：\Quote{愿上主的名永受赞颂，因为祂在祂的全体子民\Place{以色列}中彰显了你的圣洁。}\par

% structure: p0041 source=h2 target=section reason=relative-heading-rank; base=h2

\section{第十三章。}

过了不久，按照\Person{凯撒奥古斯都}的谕令进行了一次登记，命令天下人民各归本乡登记。这次登记是由\Place{叙利亚}总督\Person{基黎诺}经办的。\BibleRef{路 2:1} 因此，\Person{若瑟}应当与\Person{荣福玛利亚}一同去\Place{白冷}登记，因为他们属于\Person{犹大}支派，出自\Person{达味}的家族和家室。于是，\Person{若瑟}和\Person{荣福玛利亚}在通往\Place{白冷}的路上行走时，\Person{玛利亚}对\Person{若瑟}说：\Quote{我看见两种人民在我面前，一种哭泣，一种喜乐。}\Person{若瑟}回答说：\Quote{安静地坐在你的牲畜上，不要说多余的话。}这时，有一位身穿白衣的俊美少年出现在他们面前，对\Person{若瑟}说：\Quote{你为什么说\Person{玛利亚}关于这两种人民的话是多余的呢？因为她看见犹太人民哭泣，因为他们离开了他们的天主；而外邦人民喜乐，因为他们如今被接纳并接近了主，正如主应许我们的祖先\Person{亚巴郎}、\Person{依撒格}和\Person{雅各伯}的话：因为时候近了，\Person{亚巴郎}的后裔中，万民都要蒙福。}\BibleRef{创 12:3}\par

当他说完这些话时，天使命令牲畜停下，因为\Person{玛利亚}分娩的时候到了；天使吩咐\Person{荣福玛利亚}从牲畜上下来，进入山洞下的一个凹洞里，那里从来就没有光明，常常是黑暗，因为日光不能照到。当\Person{荣福玛利亚}进去后，那地方开始发出光芒，明亮如同白天的第六时辰。天主的光在山洞里照耀，\Person{荣福玛利亚}在那里期间，无论昼夜都不缺光明。她在那里生了一个儿子，祂出生时天使们环绕着祂。祂一出生，就站了起来，天使们朝拜祂，说：\Quote{光荣归于天主高天，平安归于地上蒙悦纳的人。}当主诞生临近时，\Person{若瑟}已出去寻找接生婆。他找到了她们，回到山洞，看见\Person{玛利亚}带着她所生的婴儿。\Person{若瑟}对\Person{荣福玛利亚}说：\Quote{我给你带来了两个接生婆——\Person{泽洛米}和\Person{撒罗默}；她们站在山洞入口外面，因着极大的光明，不敢进来。}\Person{荣福玛利亚}听了这话，微微一笑；\Person{若瑟}对她说：\Quote{不要笑，要明智地让她们来看你，免得你需要她们为你医治。}于是她吩咐她们进来。\Person{泽洛米}进来了，\Person{撒罗默}留在外面。\Person{泽洛米}对\Person{玛利亚}说：\Quote{请让我触摸你。}\Person{玛利亚}允许她检查后，那接生婆大声呼喊说：\Quote{主，上主全能者，求祢怜悯我们！从未听过也未曾想过，有人乳房充满乳汁，而所生的儿子却显示母亲仍是童贞。但祂出生时没有流血，生产时没有痛苦。童贞女怀孕，童贞女生产，童贞女依然如故。}\Person{撒罗默}听到这话，说：\Quote{请让我触摸你，证明\Person{泽洛米}说的是否真实。}\Person{荣福玛利亚}允许她触摸。当她把手从触摸中收回时，她的手枯干了，因剧痛她开始痛哭，极其痛苦，哭喊着说：\Quote{上主天主！祢知道我素来敬畏祢，我白白地照顾所有穷人；我没有从寡妇和孤儿那里拿过任何东西，也没有让有需要的人空手回去。看，我现在因我的不信而陷入悲惨，因为我无端想要试探祢的童贞。}\par

她正说着时，有一位身穿发光衣服的青年站在她身边，说：\Quote{到婴儿那里去，朝拜祂，用手触摸祂，祂会治愈你，因为祂是世界的救主，是所有仰望祂的人的救主。}她急忙走到婴儿那里，朝拜祂，触摸了包裹祂的襁褓的穗边，她的手立刻痊愈了。她出去后，开始大声呼喊，讲述她所看见和所经历的奇妙之事，以及她如何被治愈，以至于许多人因她的陈述而相信了。\par

还有一些牧人证实，他们在半夜看见天使们唱赞美诗，赞美并祝福天主，说：\BibleQuote{众人的救主已经诞生了，祂是基督主，藉着祂，救恩将归于以色列。}\BibleRef{路 2:8}\par

此外，有一颗大星，比自创世以来所见过的任何星都大，从晚上到早晨照耀在山洞上空。在\Place{耶路撒冷}的先知们说，这颗星指明了基督的诞生，祂将恢复恩许，不仅给\Place{以色列}，也给万民。\par

% structure: p0047 source=h2 target=section reason=relative-heading-rank; base=h2

\section{第十四章}

在我们主\Person{耶稣基督}诞生后的第三天，\Person{最荣福的玛利亚}走出山洞，进入一个马厩，将婴儿放在槽里，牛和驴朝拜了祂。于是应验了\Person{依撒意亚}先知所说的：\BibleQuote{牛认识主人，驴认识主人的槽。}\BibleRef{依 1:3} 因此，那些动物，牛和驴，有祂在它们中间，不停地朝拜祂。于是应验了\Person{哈巴谷}先知所说的：\BibleQuote{你在两只动物之间被显示。}\Person{若瑟}和\Person{玛利亚}在那里住了三天。\par

% structure: p0049 source=h2 target=section reason=relative-heading-rank; base=h2

\section{第十五章}

第六天，他们进了\Place{白冷}，在那里过了第七天。第八天，他们给婴儿行了割损礼，给祂起名叫耶稣，这是祂在母胎受孕前天使所称呼的名字。\BibleRef{路 2:21} \Person{玛利亚}取洁的日子满了，按照梅瑟法律，\Person{若瑟}将婴儿带到上主的圣殿。当婴儿接受了\OriginalTerm{帕里托穆斯（割损）}{parhithomus}后，他们献上了一对斑鸠或两只雏鸽。\BibleRef{肋 12:8}\par

那时，在圣殿里有一位天主的人，完善而正义，名叫西默盎，已有一百一十二岁。他从主那里得到答复，说他不得尝死味，直到他看见基督、天主子、活于肉身。并一看见这孩子，就大声呼喊说：\Quote{天主眷顾了祂的子民，主成就了祂的应许。}他急忙上前朝拜了祂。随后，祂把他抱入自己的披风，亲吻祂的脚，说：\BibleQuote{主啊，现在可照你的话，让你的仆人平安离去；因为我亲眼看见了你的救恩，即你在万民面前所预备的，为作照亮外邦人的光明，以及你的百姓以色列的荣耀。}\BibleRef{路 2:22}\par

在主的圣殿里，还有一位女先知亚纳，是法奴耳的女儿，属于阿协尔支派。她从童贞出嫁后，与丈夫同住了七年，此后守寡至今已八十四年。她从未离开过主的圣殿，而以禁食和祈祷度过时日。她也同样朝拜了这孩子，说：\Quote{在祂内有世界的救赎。}\BibleRef{路 2:36}\par

% structure: p0053 source=h2 target=section reason=relative-heading-rank; base=h2

\section{第十六章。}

第二年过后，有贤士从东方来到耶路撒冷，带来丰厚的礼物。他们向犹太人仔细查问说：\Quote{那生来作你们君王的人在何处？我们在东方看见了祂的星，特来朝拜祂。}此事传到了黑落德王那里，使他大为惊慌，于是他召集了经师、法利塞人和民间的教师，问他们先知曾预言基督应在何处诞生。他们回答说：\BibleQuote{在犹大的白冷。因为经上记载：你犹大地上的白冷，决不是犹大诸城中最小的一处，因为将由你出来一位领袖，祂将牧养我的百姓以色列。}\BibleRef{米 5:2} 于是黑落德王召见了贤士，仔细询问他们那颗星何时出现。然后，他派他们前往白冷，说：\Quote{你们去仔细寻访那孩子，一旦找到了，就回来报信，我也好去朝拜祂。}贤士们上路时，那颗星又出现在他们眼前，仿佛引导着他们，走在他们前面，直到孩子所在的地方。贤士们看见那颗星，就极其高兴；他们进了屋子，看见孩子耶稣坐在祂母亲的膝上。于是他们打开宝盒，向有福的玛利亚和若瑟献上丰厚的礼物。他们各自向孩子献上一块黄金。同样，一位献上黄金，另一位献上乳香，第三位献上没药。当他们准备回去见黑落德时，在睡梦中被一位天使警告，不要回到黑落德那里，于是他们从另一条路返回了自己的家乡。\BibleRef{玛 2:1}\par

% structure: p0055 source=h2 target=section reason=relative-heading-rank; base=h2

\section{第十七章。}

当黑落德见自己被贤士戏弄，心中怒火中烧，便派人到各条路上，想捉拿他们并处死他们。但既然找不到他们，他就再次派人到白冷及其全境，凡两岁及以下的男孩，一律按他从贤士那里查定的时间，全部杀尽。\BibleRef{玛 2:16}\par

就在此事发生的前一天，若瑟在梦中由主的使者警告，对他说：\Quote{带着玛利亚和孩子，取道旷野往埃及去。}若瑟便照天使的话去行了。\BibleRef{玛 2:14}\par

% structure: p0058 source=h2 target=section reason=relative-heading-rank; base=h2

\section{第十八章。}

他们来到一个山洞，想在那里休息，有福的玛利亚便从驴上下来，怀抱孩子耶稣坐下。若瑟身边有三个少年，玛利亚身边有一个少女，与他们同行。看，突然从洞里出来许多龙；孩子们见了，吓得大声喊叫。于是耶稣从祂母亲的怀中下来，站在龙面前；龙便朝拜耶稣，随后退去。这就应验了达味先知所说的话：\BibleQuote{地上的龙，请赞美主；龙啊，一切深渊，请赞美主。}幼童耶稣走在它们前面，命令它们不可伤害任何人。但玛利亚和若瑟非常害怕，怕孩子被龙伤害。耶稣对他们说：\Quote{不要害怕，不要以为我是个小孩子；因为我如今始终是完美的，森林中的一切野兽都必须在我面前驯服。}\par

% structure: p0060 source=h2 target=section reason=relative-heading-rank; base=h2

\section{第十九章。}

狮子与豹子同样朝拜祂，并在旷野中陪伴他们。无论若瑟和有福的玛利亚走到哪里，它们都在前面引路，低着头，摇摆尾巴表示顺服，极其虔诚地朝拜祂。起初，玛利亚看见狮子、豹子和各种野兽围拢过来，非常害怕。但婴孩耶稣以喜悦的面容注视着她，说：\Quote{母亲，不要害怕；因为它们来不是要伤害你，而是急忙要服事你我二人。}这话赶走了她心中所有的恐惧。那些狮子一直与他们同行，还有牛、驴和驮行李的牲畜，它们虽然紧挨着，却没有伤害任何一个。它们与从犹大带来的绵羊和公羊和平相处。它们走在狼群中，无所畏惧，也没有一个受其他动物伤害。这就应验了先知所说的话：\BibleQuote{豺狼将与羔羊同食，狮子和牛将一起吃草。}\BibleRef{依 65:25} 又有两头牛一同拉着一辆载有旅途用品的车，狮子为它们引路。\par

% structure: p0062 source=h2 target=section reason=relative-heading-rank; base=h2

\section{第二十章。}

行路第三天，荣福\Person{玛利亚}因沙漠中烈日过甚而疲倦；她看见一棵棕榈树，就对\Person{若瑟}说：\Quote{让我在这棵树的荫下休息片刻。}\Person{若瑟}便急忙领她到棕榈树下，扶她下了坐骑。荣福\Person{玛利亚}坐在那里，抬头看棕榈的枝叶，见满结果实，就对\Person{若瑟}说：\Quote{真希望能从这棵棕榈上摘些果子。}\Person{若瑟}对她说：\Quote{你说这话令我惊奇，你明明看见这棵棕榈有多高，竟想吃它的果子。我更担心的是缺水，因为皮囊已经空了，我们没有什么可以解渴，连牲畜也是如此。}那时，孩童耶稣，面容欢愉，躺卧在母亲的怀中，对棕榈说：\Quote{棕榈树啊，垂下你的枝条，用你的果子滋润我的母亲。}话音刚落，棕榈树就垂下树梢，直到荣福\Person{玛利亚}的脚前；他们从树上摘取果子，大家都得到了滋润。他们摘完所有果子后，棕榈树仍弯着，等待那命令它俯首者发令让它直起。耶稣便对它说：\Quote{棕榈树啊，挺起身来，要坚强，成为我父乐园中树木的同伴；并且从你的根部开启一道隐藏在地下的水脉，让水流出来，好使我们从你得到满足。}棕榈树立即挺起，它的根部开始涌出一股极其清澈、凉爽、晶莹的泉水。他们看见泉水，喜乐满怀，自己和所有牲畜都得到了满足。因此他们感谢天主。\par

% structure: p0064 source=h2 target=section reason=relative-heading-rank; base=h2

\section{第二十一章}

次日，当他们从那里启程，在开始行路的时候，耶稣转向棕榈树说：\Quote{棕榈树啊，我给你这特恩：你的一条树枝将由我的天使取下，栽种在我父的乐园中。并且我要赐你这祝福：凡在竞赛中得胜的，人将说：\InnerQuote{你已获得胜利的棕榈枝。}}祂正说着这话，看，上主的一位天使显现，站在棕榈树上；他取下一条树枝，手持树枝飞上天去。他们看见这事，就俯伏在地，如同死去。耶稣对他们说：\Quote{你们的心为什么被恐惧占据？难道你们不知道，这棵我命移植到乐园的棕榈树，将要在福乐之所为诸圣预备，如同在旷野之地为我们预备一样吗？}他们便充满喜乐，获得力量，全都站起身来。\par

% structure: p0066 source=h2 target=section reason=relative-heading-rank; base=h2

\section{第二十二章}

此后，他们行路时，\Person{若瑟}对耶稣说：\Quote{主，天气酷热；若祢愿意，我们沿着海边走，这样可以在沿海的城市休息。}耶稣对他说：\Quote{若瑟，不要害怕；我为你缩短路程，原本需三十天走完的路，你将在这一天内走完。}他们正说着话，看，他们往前看，开始看见埃及的山脉和城市。\par

他们欢欣踊跃，来到\Place{赫尔莫波利斯}地区，进入埃及一个叫\Place{索提嫩}的城市；因为在那里不认识任何人可求住宿，便进入一座称为\Place{埃及朱庇特殿}的庙宇。这庙里设立了三百五十五尊偶像，每尊偶像在其特定日子接受神性尊荣和圣礼。因为同一城市的埃及人进入朱庇特殿，祭司们告诉他们每天献上多少祭品，按照那神所受的尊敬。\par

% structure: p0069 source=h2 target=section reason=relative-heading-rank; base=h2

\section{第二十三章}

当至圣\Person{玛利亚}带着小婴孩进入庙宇时，所有偶像都俯伏在地，全都面朝下粉碎破裂（\BibleRef{撒上 5:3}），因此他们清楚显出自己毫无价值。于是应验了先知\Person{依撒意亚}所说的：\BibleQuote{看，上主必乘快速云彩来到埃及，埃及的一切偶像在祂面前战栗。}（\BibleRef{依 19:1}）\par

% structure: p0071 source=h2 target=section reason=relative-heading-rank; base=h2

\section{第二十四章}

那时，该城总督\Person{阿弗罗多修斯}听到消息，率全军来到庙宇。庙宇的祭司看见\Person{阿弗罗多修斯}率全军前来，以为他急于要看诸神因谁而跌倒下的人的报应。但当他进入庙宇，看见所有神祇都面朝下俯伏在地，便走到怀抱着主的荣福\Person{玛利亚}面前，朝拜祂，并对全军和全体朋友说：\Quote{除非这位是我们诸神的天主，我们的诸神不会在祂面前俯伏在地，也不会在祂面前这样躺卧；因此他们默默承认祂是他们的主。除非我们留心去做我们看见诸神所做的事，我们可能冒祂愤怒的危险，全都遭受毁灭，正如埃及王法老所遭遇的，他不相信这么强大的权能，连同他全军都被淹死在海中（\BibleRef{出 15:4}）。}于是那全城的人因着耶稣基督，信从了主天主。\par

% structure: p0073 source=h2 target=section reason=relative-heading-rank; base=h2

\section{第二十五章}

不久之后，天使对\Person{若瑟}说：\Quote{返回犹大地，因为那些谋取这孩子性命的人已经死了。}\par

% structure: p0075 source=h2 target=section reason=relative-heading-rank; base=h2

\section{第二十六章}

耶稣从埃及回来后，在加里肋亚，到了祂四岁那年，有一个安息日，祂与一些孩子在约但河边玩耍。祂坐在那里，用泥土为自己做了七个水池，并为每个水池做了通道，祂一命令，水就从溪流中流入水池，然后又收回。那时，有一个孩子，是魔鬼之子，因嫉妒而堵塞了供应水池的通道，推翻了耶稣所建的一切。耶稣对他说：\Quote{你有祸了，你这死亡之子，撒殚之子！你竟毁坏我所做的工程吗？}立刻，那做这事的孩子就死了。死者的父母大声哭喊，向玛利亚和若瑟抱怨，对他们说：\Quote{你的儿子咒骂了我们的儿子，他就死了。}若瑟和玛利亚听见这话，因那孩子父母的喊声和犹太人的聚集，就立刻来到耶稣那里。但若瑟私下对玛利亚说：\Quote{我不敢对祂说话；你劝诫祂，说：你为何使我们对民众怀恨？我们为何要承受人恼恨的烦扰？}祂的母亲来到祂面前，问祂说：\Quote{我主，他做了什么以致死亡？}祂说：\Quote{他该死，因为他散乱了我所做的工程。}祂的母亲又求祂说：\Quote{我主，不要这样，因为众人都起来反对我们。}但祂不愿使母亲忧伤，就用右脚踢那死孩子的臀部，对他说：\Quote{起来，你这不义之子！因为你不配进入我父的安息，因为你毁坏了我所做的工程。}那死了的人就起来，走开了。耶稣以祂权能的话语，通过水渠将水引入水池。\par

% structure: p0077 source=h2 target=section reason=relative-heading-rank; base=h2

\section{第二十七章}

此后，在众人眼前，耶稣从祂所做的水池中取泥，用泥做了十二只麻雀。那时正是安息日，耶稣做这事时，有许多孩子与祂在一起。有一个犹太人看见祂这样做，就对若瑟说：\Quote{若瑟，你看不见孩子耶稣在安息日做不许做的事吗？祂用泥土做了十二只麻雀。}若瑟听见，就责备祂说：\Quote{你为什么在安息日做我们不许做的事？}耶稣听见若瑟的话，就拍手，对麻雀说：\Quote{飞吧！}它们应声而飞。在一切站着的人眼前和耳中，祂对鸟说：\Quote{去吧，飞遍大地和整个世界，活着。}在场的人看见这样的奇事，都大为惊讶。有人赞美敬佩祂，有人辱骂祂。其中有些人去告诉司祭长和法利塞人的首领，说若瑟的儿子耶稣在以色列全体民众眼前行了伟大的征兆和奇迹。这事传遍了以色列十二支派。\par

% structure: p0079 source=h2 target=section reason=relative-heading-rank; base=h2

\section{第二十八章}

又有亚纳的儿子，圣殿的一位司祭，与若瑟同来，在众人眼前，手拿棍杖，大怒之下拆毁了耶稣亲手所做的水坝，放出了祂从溪流中积蓄的水。他还堵塞了进水的渠道，然后将它捣毁。耶稣看见，就对那毁坏祂水坝的孩子说：\Quote{你这极邪恶的不义之种！死亡之子！撒殚的工场！你的种子必无力，你的根必无湿气，你的枝子必枯干，不结果实。}立刻，在众人眼前，那孩子枯干而死。\par

% structure: p0081 source=h2 target=section reason=relative-heading-rank; base=h2

\section{第二十九章}

那时若瑟战栗，拉住耶稣，带祂回自己的家，祂的母亲也与祂同去。看，突然从对面跑来一个孩子，也是作恶的工人，撞到耶稣的肩膀上，想要戏弄祂，或伤害祂，如果可能的话。耶稣对他说：\Quote{你必不能安然无恙地从你走的路上回去。}他立刻跌倒死了。死者的父母看见所发生的事，哭喊说：\Quote{这孩子从哪里来？显然祂所说的每一句话都是真的；往往祂未说出口就已应验。}死者的父母来到若瑟那里，对他说：\Quote{带耶稣离开此地，因为祂不能与我们一起住在这城里；至少教祂祝福，不要咒骂。}若瑟来到耶稣面前，劝诫祂说：\Quote{你为什么做这样的事？如今已经有许多人因你而忧伤反对我们，我们因你而忍受人的辱骂。}耶稣回答若瑟说：\Quote{除非父亲按这时代的知教训导，没有人是智慧之子；父亲的咒骂只能伤害作恶的人。}于是众人一同起来反对耶稣，向若瑟控告祂。若瑟看见这情形，极其恐惧，害怕以色列民众的暴力和骚乱。就在那时，耶稣抓住那死孩子的耳朵，在众人眼前将他从地上提起：他们看见耶稣对他说话，如同父亲对儿子说话。他的灵魂回到他身上，他复活了。众人都惊讶。\par

% structure: p0083 source=h2 target=section reason=relative-heading-rank; base=h2

\section{第三十章}

当时，有一位名叫匝加亚斯的犹太经师听见耶稣这样说；见他因自身所具的德能而无法被驳倒，便恼怒起来，开始粗鲁、愚昧且毫无忌惮地反对若瑟。他说：\Quote{你难道不愿把你的儿子交给我，好让他受教于人文知识和敬畏之心吗？但我看出，玛利亚和你自己更关心你们的儿子，而不顾以色列民中的长老们对他所说的话。你们本该更尊敬我们——全以色列教会的长老们，好使他能与孩子们彼此相爱，并在我们中间受教于犹太学问。}若瑟却对他说：\Quote{有谁能管束这孩子、教导他呢？但若你能管束他并教导他，我们绝不阻止他跟你学习众人所学的事物。}耶稣听了匝加亚斯的话，便回答他说：\Quote{你刚才所说的律法诫命，以及你所提及的一切，都当由那些受教于人文知识的人遵守；但我是与你的法庭无干的，因为我按肉身没有父亲。你们诵读律法、精于律法，就住在律法中；但我是在律法之先的。既然你想在学识上无人能与你相比，你就要受我的教导，因为除了你所说的那些事，别人不能教导什么；但唯有那配得的人才能教导。当我在地上被高举时，我要使你们族谱的一切提及都止息。因为你们不知道自己何时出生；唯有我知道你们何时出生，以及你们在地上的寿命有多长。}凡听见这些话的人都惊讶万分，喊道：\Quote{呵！呵！呵！这真是极其伟大奇妙的奥秘！我们从未听过这样的事！从未有人说过这样的话，也从未有先知、法利塞人或经师在任何时候说过或听过。我们知道他的来历，他才不过五岁，怎能说出这些话呢？}法利塞人回答说：\Quote{我们从未听过任何一个如此年幼的孩子说出这样的话。}耶稣回答他们说：\Quote{你们因一个孩子说出这样的话就感到惊奇吗？那么，你们为什么不信我对你们所说的话呢？你们大家都因我对你们说\InnerQuote{我知道你们何时出生}而惊讶。我要告诉你们更大的事，使你们更加惊奇：我见过亚伯拉罕，就是你们称为你们祖先的那位，并与他谈过话；他也见过我。}\BibleRef{若 8:56}他们听了这话，就都默不作声，没有一个人敢说话。耶稣对他们说：\Quote{我曾在孩子们中间与你们同在，你们却不认识我；我向你们说话，如同对智慧人说话，你们却不明白我的话；因为你们比我还年轻，且信心小。}\par

% structure: p0085 source=h2 target=section reason=relative-heading-rank; base=h2

\section{第三十一章}

老师宰喀雅斯，法律博士，再次对若瑟和玛利亚说：\Quote{把孩子交给我，我将把他交给老师肋未，由他教他字母和教导他。}于是若瑟和玛利亚安抚耶稣，带祂到学校，让年老的肋未教祂字母。耶稣一进去，便沉默不语。老师肋未向耶稣念了一个字母，从第一个字母阿肋弗开始，对祂说：\Quote{回答。}但耶稣沉默，一言不答。因此教师肋未发怒，拿起他的一根檉柳棍，击打祂的头。耶稣对老师肋未说：\Quote{你为什么打我？你当知道，被打的人更能教导打他的人，远胜于被他教导。因为我能教导你所说的那些事。但所有这些说话与听话的人都是瞎子，如同响亮的铜钹或铿锵的铙钹，其中没有对所发声音意义的知觉。}耶稣又对宰喀雅斯说：\Quote{从阿肋弗直到特特，每个字母都是按顺序知道的。所以请先说特特是什么，我就告诉你阿肋弗是什么。}耶稣又对他们说：\Quote{那不知道阿肋弗的人，怎能说特特呢？你们这些伪善者！告诉我第一个字母阿肋弗是什么；等你们说出了贝特，我就相信你们。}耶稣开始逐一询问字母的名称，说：\Quote{让法律博士告诉我们第一个字母是什么，或者它为什么有许多三角形、分级、近锐、中间、外凸、伸长、直立、横卧、弯曲。}肋未听了，对这些字母名称的排列惊得目瞪口呆。然后他在众人面前开始喊叫说：\Quote{这样的人怎能在世上活着？他实在该被挂在十字架上。因为祂能灭火，并戏弄其他刑罚方式。我想祂在大洪水之前就已存在，在洪水之前就已出生。是哪个子宫怀了祂？哪个母亲生了祂？哪个乳房给祂哺乳？我在祂面前逃跑，我受不了祂口中所说的话；我的心听到这些话就惊惶。我不认为任何人能理解祂所说的，除非天主与他同在。我这可怜的人，现在成了祂的笑柄。因为当我以为有了一个学生时，我却不认识祂，反而找到了我的老师。我还能说什么？我抵挡不了这孩子的言语；我现在要逃离这城，因为我无法理解它们。像我这样的老人竟被一个孩童打败，因为我找不到祂所说的话的开头或结尾。因为要找到祂自己的开头并非易事。我确实告诉你们，我没有说谎：在我看来，这孩子的行为、祂谈话的开头和祂意图的结局，与凡人毫无共同之处。因此我不知道祂是巫师还是神，或者至少是天使在祂内说话。祂是谁，从何而来，将来会成为什么样的人，我都不知道。}于是耶稣面带喜色，微笑着用命令的口气对站在旁边听的所有以色列子民说：\Quote{愿不结果子的结果子，盲人看见，瘸子端正行走，穷人享受今世的美物，死人复活，使每个人回到他原来的状态，并住在祂内，祂是生命和永恒甘饴之根。}孩童耶稣说了这话，所有患有恶疾的人立刻痊愈了。他们再不敢对祂说什么，也不敢再听祂的话。\par

% structure: p0087 source=h2 target=section reason=relative-heading-rank; base=h2

\section{第三十二章}

此后，若瑟和玛利亚带着耶稣离开那里，来到纳匝肋城；祂与父母留在那里。在七日的第一日，当耶稣与一些孩子在某个屋顶上玩耍时，其中一个孩子将另一个孩子从屋顶推下地面，那孩子死了。死者的父母没有看见这事，便向若瑟和玛利亚喊叫说：\Quote{你们的儿子把我们儿子推到地上，他死了。}但耶稣沉默，一言不答。若瑟和玛利亚急忙来到耶稣面前；祂的母亲问祂说：\Quote{我的主，告诉我你是否把他推下去了。}耶稣立刻从屋顶下到地面，叫那孩子的名字：则诺。那孩子回答说：\Quote{我的主。}耶稣对他说：\Quote{是我把你从屋顶推到地上的吗？}他说：\Quote{不是，我的主。}那已死孩子的父母惊奇，因所行的奇迹而尊荣耶稣。若瑟和玛利亚便带着耶稣离开那里，往耶里哥去了。\par

% structure: p0089 source=h2 target=section reason=relative-heading-rank; base=h2

\section{第三十三章}

那时耶稣六岁，祂母亲打发祂拿着水罐与孩子们到井边打水。祂打水后，一个孩子冲撞祂，打碎了水罐。但耶稣展开祂所穿的外衣，在外衣里装了原先水罐里那么多的水，拿回给祂母亲。她看见便惊奇，心中思忖，把这一切都记在心里。\BibleRef{路 2:19}\par

% structure: p0091 source=h2 target=section reason=relative-heading-rank; base=h2

\section{第三十四章}

又有一天，祂到田里去，从祂母亲的粮仓里拿了一点麦子，自己撒种。麦子发芽、生长，并大大增加。最后，祂亲自收割，收获了三个\Quote{科尔}，分给祂的许多熟人。\par

% structure: p0093 source=h2 target=section reason=relative-heading-rank; base=h2

\section{第三十五章}

有一条路从耶里哥出来，通向约但河，就是以色列子民渡过的地方；据说约柜曾在那里安放。耶稣那时八岁，祂从耶里哥出来，往约但河去。路边靠近约但河岸有一个岩洞，里面有一只母狮正在哺育幼崽；没有人敢安全地走那条路。耶稣从耶里哥来，知道那岩洞里有母狮生了幼崽，就在众人眼前进入洞中。狮子们看见耶稣，就跑到祂面前，朝拜祂。耶稣坐在洞里，狮子的幼崽在祂脚边跑来跑去，向祂表示亲热，并嬉戏。那些较大的狮子低着头，站在远处，朝拜祂，并摇尾向祂表示亲热。那时，站在远处的人看不见耶稣，就说：除非祂或祂的父母犯了重罪，祂不会自愿把自己交给狮子。当人们这样心里思量，并陷于极大忧愁时，忽然在众人眼前，耶稣从洞中出来，狮子走在祂前面，狮子的幼崽在祂脚前互相嬉戏。耶稣的父母低头站在远处观看；同样，人们也因狮子而站在远处，因为他们不敢靠近它们。于是耶稣开始对众人说：禽兽比你们好得多，因为它们认出它们的主，并光荣祂；而你们人，按天主的肖像和模样被造，却不认识祂！野兽认识我，且变得驯良；人看见我，却不承认我。\par

% structure: p0095 source=h2 target=section reason=relative-heading-rank; base=h2

\section{第三十六章}

这些事以后，耶稣在众人眼前与狮子一同渡过约但河；约但河的水左右分开。然后耶稣对狮子说，在众人听见下：平安去吧，不要伤害任何人；但人也不可伤害你们，直到你们回到你们出来的地方。狮子们不仅用动作，也用声音向祂告别，然后回到自己的地方。耶稣却回到祂母亲那里。\par

% structure: p0097 source=h2 target=section reason=relative-heading-rank; base=h2

\section{第三十七章}

若瑟是木匠，他用木头只做牛轭、犁、农具和木床。有一个青年吩咐他做一张六肘长的床榻。若瑟命令他的仆人用铁锯按着送来的尺寸锯木头。但仆人没有遵守指定的尺寸，反而把一块木头锯得比另一块短。若瑟感到困惑，开始考虑怎么办。耶稣看见他这样沉思，知道这事对他是无法做到的，就用安慰的话对他说：来，我们握住木头的两端，把两端对起来，使它们正好吻合，然后拉向我们，因为我们能使它们一样长。于是若瑟照吩咐做了，因为他知道祂能做成祂所愿意的任何事。若瑟握住木头的两端，把它们靠墙对着自己；耶稣握住木头的另一端，把较短的那块拉向祂，使它和较长的那块一样长。祂对若瑟说：去工作吧，做你承诺要做的事。若瑟就做了他所承诺的事。\par

% structure: p0099 source=h2 target=section reason=relative-heading-rank; base=h2

\section{第三十八章}

第二次，若瑟和玛利亚被众人请求，让耶稣在学校里学习识字。他们没有拒绝；按照长老们的吩咐，他们把祂带到一位老师那里去接受人的学问。老师用命令的口吻开始教导祂说：说\Quote{阿尔法}。耶稣对他说：你先告诉我什么是\Quote{贝塔}，我就告诉你什么是\Quote{阿尔法}。老师因此发怒，打了耶稣；他一打祂，就倒地死了。耶稣又回家到祂母亲那里。若瑟害怕，就把玛利亚叫到跟前，对她说：你确实知道，我为这孩子忧愁以至死。因为很可能什么时候有人恶意打祂，祂就会死。但玛利亚回答说：天主的人啊！不要相信这是可能的。你确实可以相信，那派遣祂生于人间的那一位，祂自己会保护祂免于一切祸害，并因祂自己的名保护祂免于邪恶。\par

耶稣又回家到祂母亲那里。若瑟害怕，就把玛利亚叫到跟前，对她说：你确实知道，我为这孩子忧愁以至死。因为很可能什么时候有人恶意打祂，祂就会死。但玛利亚回答说：天主的人啊！不要相信这是可能的。你确实可以相信，那派遣祂生于人间的那一位，祂自己会保护祂免于一切祸害，并因祂自己的名保护祂免于邪恶。\par

% structure: p0102 source=h2 target=section reason=relative-heading-rank; base=h2

\section{第三十九章}

第三次，犹太人又请求玛利亚和若瑟哄劝祂去另一位老师那里学习。若瑟和玛利亚因惧怕百姓、首领们的蛮横和司铎们的威胁，就又把祂带到学校，知道祂不能从人那里学到什么，因为祂只有从天主才有完备的知识。耶稣受圣神引导进入学校后，祂从那位教授法律的老师手中拿过书，在众人眼前和耳边开始读，不是读他们书中写的，而是因永活天主之神发言，如同活泉涌出的水流，而泉源常满。祂用这样的能力教导百姓永活天主的伟大事物，以至于老师本人俯伏在地朝拜祂。坐着听祂说这些话的百姓，心中转为惊奇。若瑟听见这事，就跑来见耶稣，怕老师本人死了。老师看见若瑟，对他说：你给我的不是学生，而是老师；谁能抵挡祂的话呢？于是应验了圣咏作者所说的话：\Quote{天主的河满了水；祢为他们预备了五谷，因为如此是它的供应。}\par

% structure: p0104 source=h2 target=section reason=relative-heading-rank; base=h2

\section{第四十章}

此后，\Person{若瑟}带着\Person{玛利亚}和\Person{耶稣}离开那里，前往海边的\Place{葛法翁}，因为敌人的恶意。\Person{耶稣}住在\Place{葛法翁}时，城内有一个名叫\Person{若瑟}的人，极其富有。他却因患病而消瘦，死了，躺在自己的床上。\Person{耶稣}听见城中有人为死者哀悼、哭泣、悲叹，便对\Person{若瑟}说：\Quote{你为什么不肯将你恩惠的好处给予这人，既然他和你同名？}\Person{若瑟}回答祂说：\Quote{我有什么能力或本事能给他好处呢？}\Person{耶稣}对他说：\Quote{拿起你头上的手帕，去放在死者的脸上，对他说：\InnerQuote{基督治愈你。}那死者便会立刻痊愈，从床上起来。}\Person{若瑟}听了这话，便遵\Person{耶稣}的命令去了，跑进死者的家，将他戴在头上的手帕放在躺在床上的人的脸上，说：\Quote{耶稣治愈你。}那死者立刻从床上起来，问\Person{耶稣}是谁。\par

% structure: p0106 source=h2 target=section reason=relative-heading-rank; base=h2

\section{第四十一章}

他们离开\Place{葛法翁}，来到名叫\Place{白冷}的城；\Person{若瑟}和\Person{玛利亚}住在他自己的房子里，\Person{耶稣}也和他们同住。有一天，\Person{若瑟}叫来他的长子\Person{雅各伯}，派他到菜园里去摘菜，好做汤。\Person{耶稣}跟着祂的兄弟\Person{雅各伯}进了菜园，\Person{若瑟}和\Person{玛利亚}却不知道。\Person{雅各伯}正在摘菜时，一条毒蛇突然从洞里出来，咬了他的手，\EditorialNote{\BibleRef{宗 28}}，他因剧痛开始喊叫。他精疲力竭，痛苦地喊着说：\Quote{哎呀！哎呀！一条该死的毒蛇咬了我的手。}\Person{耶稣}正站在他对面，听到痛苦的喊声，就跑向\Person{雅各伯}，抓住他的手；祂所做的只是向\Person{雅各伯}的手吹气，使它凉下来：\Person{雅各伯}立刻痊愈了，蛇也死了。\Person{若瑟}和\Person{玛利亚}不知道所发生的事；但听到\Person{雅各伯}的喊声和\Person{耶稣}的命令，他们跑到菜园，发现蛇已经死了，\Person{雅各伯}完全痊愈了。\par

% structure: p0108 source=h2 target=section reason=relative-heading-rank; base=h2

\section{第四十二章}

\Person{若瑟}与他的儿子们——\Person{雅各伯}、\Person{若瑟}、\Person{犹达}、\Person{西默盎}和他的两个女儿——来赴宴时，\Person{耶稣}与祂的母亲\Person{玛利亚}，以及她的姊妹——克罗帕的玛利亚——遇到了他们；主天主曾将克罗帕的玛利亚赐给她的父亲\Person{克罗帕}和母亲\Person{亚纳}，因为他们曾将\Person{耶稣}的母亲\Person{玛利亚}奉献给主。她取了同样的名字\Person{玛利亚}，为了安慰她的父母。他们聚集后，\Person{耶稣}圣化并祝福了他们，祂是第一个开始吃喝的；因为没有人敢吃喝、入席或擘饼，直到祂圣化他们并首先这样做。如果祂碰巧不在，他们就等候，直到祂这样做。当祂不愿意来用膳时，\Person{若瑟}、\Person{玛利亚}和祂的兄弟（\Person{若瑟}的儿子们）也都不来。事实上，这些兄弟把祂的生活当作明灯放在眼前，观察祂，敬畏祂。\Person{耶稣}睡觉时，无论白天黑夜，天主的光辉照在祂身上。愿一切赞美和荣耀归于祂，直到永远。\OriginalTerm{阿们}{Amen}，\OriginalTerm{阿们}{Amen}。\par

\SourceRecord{Translated by Alexander Walker.；Ante-Nicene Fathers；Vol. 8.；Edited by Alexander Roberts, James Donaldson, and A. Cleveland Coxe.；Buffalo, NY: Christian Literature Publishing Co.,；1886.；<http://www.newadvent.org/fathers/0848.htm>.}
